\documentclass[11pt]{letter}
\usepackage[margin=1in]{geometry}
\usepackage[hidelinks]{hyperref}
\signature{Liang Zhang\University of Arizona\liangzhang1@arizona.edu}
\address{Liang Zhang\University of Arizona\Tucson, AZ, USA}
\begin{document}
\begin{letter}{Editors-in-Chief\Energy and Buildings}
\opening{Dear Editors,}

I am pleased to submit the manuscript entitled \emph{``Virtual BAS: a physics-grounded, real-world-like building
automation system emulator with labelled faults, sensor imperfections and stochastic operation for developing and
testing building AI''} for consideration as a research article in \emph{Energy and Buildings}.

Artificial intelligence for building operation---fault detection and diagnosis, supervisory and reinforcement-learning
control, and emerging ambient and large-language-model agents---must work on the data of real building automation
systems, where sensors drift, data go missing, occupants behave randomly and equipment fails. Real BAS data rarely come
with ground truth, while existing simulation test beds are mostly idealised. The manuscript presents an open-source
framework that couples a step-wise EnergyPlus simulation (the physics) with an explicitly modelled and fully labelled
real-world layer, in which BAS controllers act on imperfect observations so that the physical consequences of
imperfection are simulated. Faults and events from a library of 44 mathematically defined types can be injected at run
time by people and by AI agents, the clock runs from real time to tens of thousands of times faster, and any EnergyPlus
model is converted automatically. Beyond the points of a typical BAS, the emulator exposes the complete indoor and
outdoor environment---comfort indices, particles, VOC, light, sound, occupant counts and feedback, an extended weather
station with forecasts, sub-meters and grid signals---with every point tagged as standard, advanced or
next-generation instrumentation. Using the DOE reference large office, we verify the emulator against native
EnergyPlus and demonstrate equipment-level fault impacts, labelled data generation for FDD, AI agents in the loop, indoor air quality under wildfire smoke,
time-acceleration accuracy and generality across five reference buildings.

We believe the work fits the scope of \emph{Energy and Buildings} on building energy systems, controls, automation and
data-driven methods, and that the open software and tutorials
(\url{https://github.com/TensorBuildLab/Virtual-BAS}) will be useful to the journal's readership.

The manuscript is original, has not been published previously and is not under consideration elsewhere. The author
declares no competing interests. The use of generative AI assistance is disclosed in the manuscript.

\closing{Sincerely,}
\end{letter}
\end{document}
